% 第 8 章 核磁共振方法 —— 分片 c
% 覆盖：S1 §8.6（书内 p.555--566 / PDF p.47--58）
% 分片 a 已写章首（chapterstart）；本分片直接从 \subsection 开始，不重复章首。

\subsection{MAS NMR 技术研究积炭引起的分子筛失活}
\src{S1 §8.6，书内 p.555--562 / PDF p.47--54}

催化工业对催化剂的要求不只是活性与选择性，还包括高稳定性与长寿命，因此催化剂失活是一个十分重要的技术和经济问题。在加氢裂化和甲醇转化等催化过程中，分子筛催化剂的积炭往往是降低催化剂活性进而使其失活的主要因素。积炭指反应中生成的各种有机物沉积在催化剂表面从而引起催化剂失活；积炭的形成及其对分子筛的作用，取决于分子筛的孔结构、反应物的性质和反应条件（温度、压力等）。

失活机理并不唯一，两个对照例子足以说明。氢型丝光沸石的快速失活来自孔道内积炭分子的存在：所有活性位都不可能再与反应物接触，这是孔堵塞（pore blocking）引起的失活。HY 与 HZSM-5 分子筛虽然都有三维晶体结构，积炭的性质以及积炭与失活的速率却完全不同：HY 的积炭是非均匀的，最强的酸位首先受积炭影响；HZSM-5 的活性位强度相等，积炭是均匀的。这两种分子筛的失活不能用孔堵塞解释，而是活性位中毒（active-site poisoning）。

于是研究任务可以拆成四个问题：积炭\textbf{是什么}（碳的化学环境）、\textbf{有多少}（量）、\textbf{积在哪}（孔道内还是外表面、窗口还是笼内）、\textbf{以什么方式}让催化剂失活（酸位中毒还是孔堵塞）。多种化学和物理技术都能回答其中一部分，而固体高分辨 NMR 是最重要的方法之一：\ce{^{13}C}、\ce{^{27}Al}、\ce{^{29}Si}、\ce{^{1}H}、\ce{^{129}Xe} MAS NMR 均可用来研究积炭引起的催化剂失活。

\begin{probout}
\textbf{问题}：同一份积炭样品，为什么要换四种核去测？每一种核各自能回答上面四个问题中的哪一个，它的答案又凭什么可信？

\textbf{输出}：一张「核种--可比量--定量前提--主要干扰」的对照表（表~\ref{tab:nuclei-coke}），以及一条判读纪律：任何关于「积炭量」「酸位数」「孔容」的数字，都必须同时报出它的核种、定量前提与未消除的干扰；只给数字不给前提的结论不可用。
\end{probout}

四种核之所以不能互相替代，是因为各自只对积炭问题的一个面敏感。\ce{^{13}C} 直接看积炭本身的碳骨架，给出脂肪碳／烯碳／芳香碳／聚芳烃的分布，但它几乎不回答「积在哪」，而且积炭越重信号丢得越多；\ce{^{29}Si} 与 \ce{^{27}Al} 看的是骨架这一侧被扰动成什么样子，反映骨架弛豫的变化、脱铝以及骨架铝配位对称性的下降；\ce{^{1}H} 看氢，给出积炭后仍存活的 Brönsted 酸位数量、H/C 值，以及积炭在孔道内还是外表面；\ce{^{129}Xe} 自己不参与反应、不改变酸性，只作为惰性探针报告孔道空间被占据的几何后果：窗口是否被堵、笼的有效直径是否缩到与 Xe 原子相当、是否在积炭内部生成了新的微孔或中孔。四者的可观测量不同、定量前提不同、失效方式也不同，所以「四种核」不是四次重复测量，而是四个互补的观测面。

\subsubsection{\ce{^{13}C} MAS NMR 研究分子筛积炭}
\src{S1 §8.6.1，书内 p.555--556 / PDF p.47--48}

\begin{principle}
\textbf{核性质与观测窗口。} \ce{^{13}C} 核自旋 $I=1/2$，无四极矩，固体中只受偶极耦合、化学位移各向异性与弛豫影响，MAS 下可获得高分辨谱。天然丰度 1.10\%，检测灵敏度约为质子的 \num{1.59e-2}，因此分子筛积炭研究多用 \ce{^{13}C} CP/MAS NMR 技术并尽量提高磁场。\ce{^{13}C} 化学位移范围为 0--300 ppm，由谱峰位置区分积炭的类型，如脂肪碳、烯碳、芳香碳或聚芳烃（表~\ref{tab:c13-assign}）。
\tocheck{未对照原页面图像：灵敏度的指数（$1.59\times10^{-2}$）按 OCR 上下文与常规值重建，其参考基准（等磁场还是等频率、是否已含天然丰度因子）原书未说明}

\textbf{交叉极化的信号从哪来，为什么它不能定量。} CP 把 \ce{^{1}H} 的较大极化经偶极耦合转移给 \ce{^{13}C}，灵敏度可提高数倍，理论上限约为 $\gamma_{\mathrm{H}}/\gamma_{\mathrm{C}}\approx 4$；转移在旋转坐标系中两核的 Rabi 频率相等时最有效，即 Hartmann--Hahn 条件 $\gamma_{\mathrm{H}}B_{1\mathrm{H}}=\gamma_{\mathrm{C}}B_{1\mathrm{C}}$。因此接触时间需要优化，而转移效率还取决于 \ce{^{1}H}--\ce{^{13}C} 偶极耦合的强度与两核的弛豫行为。这就是「CP 谱峰高不等于碳原子数之比」的根源：甲基、季碳、芳香碳，以及与石墨或顺磁中心相邻的碳，极化转移效率各不相同。
\tocheck{未对照原页面图像：Hartmann--Hahn 条件与 $\gamma_{\mathrm{H}}/\gamma_{\mathrm{C}}\approx 4$ 按标准形式补入，原书本节只以文字说明「不同类型碳的交叉极化效果不同」}
\end{principle}

积炭量增加时脂肪碳峰与芳香碳峰的消长，是 \ce{^{13}C} 谱上最直观的读数。图~\ref{fig:c13-coke-model} 用洛伦兹线型叠加的模型计算给出这一趋势，同时给出被 CP 效率与信号损失压低的整体强度。

\begin{figure}[tbp]
\centering
\begin{tikzpicture}
% declare function 必须放在 \begin{tikzpicture} 之后、\begin{axis} 之前，用 \tikzset 声明；
% 线型用有界洛伦兹 1/(1+((x-a)/w)^2)，避免小 w 的高斯指数触发 Dimension too large。
\tikzset{declare function={
  pk(\x,\a,\b) = 1/(1+((\x-\a)/\b)^2);
  spa(\x) = 1.00*pk(\x,25,11) + 0.12*pk(\x,140,16);
  spb(\x) = 0.82*pk(\x,25,12) + 0.40*pk(\x,138,18);
  spc(\x) = 0.55*pk(\x,26,13) + 0.64*pk(\x,135,22);
  spd(\x) = 0.20*pk(\x,30,15) + 0.36*pk(\x,133,28);
}}
\begin{axis}[
  width=\onecolwidth, height=50mm,
  xmin=0, xmax=300, ymin=0, ymax=1.25,
  x dir=reverse,
  xlabel={\ce{^{13}C} 化学位移 / ppm},
  ylabel={相对强度 / 任意单位},
  xtick distance=50, ytick={0,0.5,1.0},
  tick label style={font={\fontsize{8}{9.6}\selectfont}},
  label style={font={\fontsize{8}{9.6}\selectfont}},
  legend style={at={(0.02,0.03)}, anchor=south west, draw=none,
     fill=white, fill opacity=0.9, text opacity=1,
     font={\fontsize{8}{9.6}\selectfont}, row sep=0.5pt, inner sep=1.5pt},
  every axis plot/.append style={thick, no markers, samples=200, domain=0:300},
]
\addplot[black] {spa(x)};
\addlegendentry{积炭 0.85\%}
\addplot[black!70, dashed] {spb(x)};
\addlegendentry{积炭 9.6\%}
\addplot[black!45, dotted] {spc(x)};
\addlegendentry{积炭 16.5\%}
\addplot[black!30, dashdotted] {spd(x)};
\addlegendentry{积炭 23.6\%}
% 承载文字的 node 一律给 text width，否则最长一行会决定整图宽度。
\node[font={\fontsize{8}{9.6}\selectfont}, text width=16mm, align=center, anchor=south]
  at (axis cs:52,1.02) {脂肪碳\\10--40 ppm};
\draw[->, thin] (axis cs:196,0.82) -- (axis cs:145,0.66);
\node[font={\fontsize{8}{9.6}\selectfont}, text width=20mm, align=left, anchor=west]
  at (axis cs:200,0.84) {芳香碳\\125--170 ppm};
\end{axis}
\end{tikzpicture}
\caption{模型计算的 \ce{^{13}C} MAS NMR 谱随积炭量的变化示意\modelmark。读者应看两件事：右侧脂肪碳峰（10--40 ppm）相对下降，
而左侧芳香碳峰（125--170 ppm）从接近本底长成主峰，说明积炭骨架随积炭量增加而芳构化；同时全谱总强度按可检出比例衰减
（积炭 23.6\% 时只有约 29\% 的积炭可由 NMR 检出）。曲线为洛伦兹线型叠加的模型计算，峰位与半峰宽按本节给出的实验区间设定；
峰高之比的绝对值不代表真实碳量之比（受 CP 效率差异控制）。\figsrc{参数取值参考 S1 §8.6.1，书内 p.555--556 / PDF p.47--48}}
\label{fig:c13-coke-model}
\end{figure}

\begin{table}[htbp]
\centering
\caption{\ce{^{13}C} 化学位移与积炭归属（本节实例）。读者应看两件事：不同类型的碳落在不同位移区间，
而「芳香碳区很宽」本身就是聚合芳烃生成的判据，比单个峰的位置更有信息量。}
\label{tab:c13-assign}
{\srcsize
\begin{tabularx}{\columnwidth}{@{}lX@{}}
\toprule
化学位移 & 归属（实例） \\
\midrule
13.5 / 24.0 / 33.0 ppm & 乙烯在脱水 HZSM-5 上 295 K 吸附生成的直链烷烃聚合物（谱峰很宽，是聚合物的谱线特征） \\
10--40 ppm & HZSM-5 上甲醇转化积炭的脂肪族化合物（分布较宽，支链烷烃多于直链烷烃） \\
13--25 ppm & H 型丝光沸石上分布较窄的脂肪烃 \\
约 150 ppm & 直链与支链烯烃 \\
125--145 ppm & HZSM-5 上的芳香化合物 \\
130--170 ppm & H 型丝光沸石的宽芳香碳区，表明存在聚合芳烃 \\
\bottomrule
\end{tabularx}}
\end{table}

\begin{measure}
\textbf{脉冲与条件}：原位池或密封转子；CP/MAS 配高功率 \ce{^{1}H} 去偶与尽可能高的转速，把旋转边带移出关注区（必要时加 TOSS 抑制边带）。

\textbf{必须记录的量}：峰位（归属）、半峰宽、边带间隔（等于转速）与边带相对强度、以及 CP 谱与直接极化谱的强度比。缺转速记录的谱无法判断哪几个峰是边带。

\textbf{如果目标是碳量而不是「有没有」}：必须另做直接极化（单脉冲）谱，循环延迟按最慢的 \ce{^{13}C} $T_1$ 取（芳构化、石墨化程度高的样品常需数十秒以上），并用已知碳量的标样，或与热重、元素分析结果对照标定。只报 CP 谱的峰面积比，不能给出积炭量。
\end{measure}

\begin{dataflow}
\textbf{（1）孔结构决定积炭类型。} 1982 年 Derouane 等首先用原位 \ce{^{13}C} MAS NMR 确定甲醇转化过程中 HZSM-5 与 H 型丝光沸石上的积炭。HZSM-5 上的积炭是分布较宽的脂肪族化合物（10--40 ppm）、一些芳香化合物（125--145 ppm）以及直链与支链烯烃（约 150 ppm），其中支链烷烃多于直链烷烃；H 型丝光沸石上的积炭是分布较窄的脂肪烃（13--25 ppm），但芳香碳区很宽（130--170 ppm），表明存在聚合芳烃。同样的反应物，在不同孔结构里给出了不同的碳骨架。

\textbf{（2）预处理改变同一分子筛上的积炭。} 乙烯在 HZSM-5 上 295 K 吸附并生成直链烷烃聚合物（13.5、24.0、33.0 ppm），宽谱峰显示聚合物的谱线特征；分子筛吸水后经 573 K 处理，谱峰明显窄化，说明高分子量聚合物已裂解生成直链与支链低碳脂肪烷烃，同时还观察到乙苯共振峰。

\textbf{（3）温度改变积炭走向，铝含量只影响其中一支。} 乙烯在 H 型丝光沸石上原位反应：500 K 以下的低温积炭多为直链与支链烷烃，500 K 以上的高温积炭多为聚合芳烃；低温积炭不受分子筛铝含量影响，高温积炭与铝含量有关。可见积炭类型是温度与酸位密度的函数，只报一个「积炭量」不足以描述样品。

\textbf{（4）可检出比例随积炭量下降，损失有三种机制。} HZSM-5 上甲醇转化：积炭量为 0.85\% 时可检测到所有的积炭；积炭量达到 23.6\% 的最高值时，只有 29\% 的积炭可由 NMR 检测。损失的原因是多方面的：其一，\ce{^{1}H} 的自旋--自旋弛豫时间太短，降低了 \ce{^{1}H}$\rightarrow$\ce{^{13}C} 交叉极化的效率；其二，高积炭量时在分子筛外表面生成导电性的石墨结构，干扰了探头的调谐，降低了 \ce{^{13}C} 检测灵敏度；其三，高积炭量的结构中含有芳环离域电子与缺陷位的顺磁中心，降低了 \ce{^{1}H}$\rightarrow$\ce{^{13}C} 交叉极化效率。能够由 NMR 检测的积炭基本上是在分子筛孔道内的；高含量积炭的分子筛表面物种主要是石墨，因此没有 NMR 信号，而这些石墨堵塞了分子筛表面的孔口，阻碍甲醇在通道内的扩散和产物扩散出孔道，从而使分子筛失活。不同类型分子筛上积炭的 \ce{^{13}C} 信号损失程度也不同。

\textbf{（5）结论的正确表述。} \ce{^{13}C} NMR 提供的信息基本上是积炭的化学特征，它取决于分子筛的孔结构、反应物的性质和反应温度；在某些情况下还可以提供积炭中间物的信息。它的局限同样明确：难以得到准确的积炭量数据，只有在轻微积炭或低温积炭的情况下，才可以粗略估计积炭量并确定积炭性质，而且不能区分分子筛孔道内和外表面积的炭。
\end{dataflow}

\begin{pitfall}
\textbf{（1）CP 谱不能直接定量。} 不同类型碳的交叉极化效果不同，\ce{^{13}C} CP/MAS NMR 谱无法用于定量确定积炭量。要把「脂肪碳峰下降、芳香碳峰上升」读成「脂肪碳减少、芳香碳增加」，必须假定各组分的 CP 效率不变；当积炭本身改变了弛豫与顺磁环境时，这个假定并不成立。定量结论应来自直接极化谱，或来自热重、元素分析等独立方法。

\textbf{（2）旋转边带会被误当成真实碳物种。} 边带以转速为间隔出现在真实峰两侧，强度随转速与 CP 条件改变，位置随转速移动；真实碳物种的峰位不随转速移动。再现聚合芳烃、烯碳这类「新物种」之前，先用两个转速对照排除边带。

\textbf{（3）「NMR 不可见」不等于「没有积炭」。} 高积炭量时外表面石墨与顺磁中心使信号消失，可检出量只是积炭的一个子集。特别要注意：23.6\% 积炭只有 29\% 可检出，但本书明确指出可检出的积炭基本位于孔道内，说明「看不见」的那部分并不都等于外表面石墨，其中一部分是孔道内碳因 CP 效率下降而丢失的信号；把差额直接当作外表面积炭会高估石墨量。

\textbf{（4）探头调谐变化属于仪器层面。} 导电石墨改变探头调谐，从而改变灵敏度；跨样品比较峰强前必须重新调谐，并记录调谐与接收增益条件，否则测到的「积炭效应」里混有仪器效应。
\end{pitfall}

\subsubsection{\ce{^{29}Si} MAS NMR 与 \ce{^{27}Al} MAS NMR 研究分子筛积炭}
\src{S1 §8.6.2，书内 p.556--558 / PDF p.48--50}

填充了孔道的积炭影响骨架 \ce{^{29}Si} 和 \ce{^{27}Al} 核的周围环境，使它们的 MAS NMR 谱发生变化，因此这两种技术可用来研究积炭，提供的是\textbf{积炭对分子筛骨架的影响}，而不是积炭本身的化学信息。

\paragraph{（1）\ce{^{29}Si} MAS NMR：骨架弛豫的变化与 H/C}

\begin{principle}
\ce{^{29}Si} 自旋 $I=1/2$，化学位移由骨架硅的 Q$^n$ 环境（铝邻居数）决定，因此谱形本身就是骨架连通性的读数。积炭对它的影响有两条彼此独立的通道：一是孔道内物种改变骨架硅周围的磁环境，使峰位与峰形变化；二是积炭取代了孔道中的顺磁性氧，降低了 \ce{^{29}Si} 和 \ce{^{1}H} 的自旋--晶格弛豫速率，增大了 Si 的自旋--晶格弛豫时间 $T_1$，于是在一定的脉冲循环时间（如 \qty{10}{s}）内磁化矢量不能充分弛豫，从而降低了信号强度。第二条通道提醒我们：信号下降未必是硅少了。

\ce{^{1}H}$\rightarrow$\ce{^{29}Si} 交叉极化（\ce{^{29}Si} CP/MAS NMR）信号的增强与吸附物中氢的含量有关，这一结果可以扩展到积炭研究：CP 信号强度对应于每个晶胞中含氢原子的数目，与积炭量（碳的数目）相比即可得到 H/C 值：
\[
\frac{\mathrm{H}}{\mathrm{C}}=\frac{N_{\mathrm{H}}}{N_{\mathrm{C}}}
\]
其中 $N_{\mathrm{H}}$ 为每个晶胞中的氢原子数（由 \ce{^{29}Si} CP 信号强度标定），$N_{\mathrm{C}}$ 为同一晶胞中积炭碳原子数（由积炭质量分数与晶胞组成换算）。
\tocheck{未对照原页面图像：H/C 的换算式、所用晶胞组成与标定参考样品原书未给出，此处按定义重建；\ce{^{29}Si} 天然丰度与化学位移范围亦未在 §8.6.2 重复给出}
\end{principle}

\begin{dataflow}
\textbf{（1）谱形的变化有阈值。} HZSM-5 分子筛积炭前后的 \ce{^{29}Si} MAS NMR 谱（图 8-17）显示：积炭在 10\%（质量分数）以下时 \ce{^{29}Si} 谱变化不大；积炭到 16.5\% 和 23.6\% 时，谱峰移向低场，峰形发生扭曲并且谱线增宽，半峰宽由 \qty{143}{Hz} 增加到 \qty{218}{Hz}，峰形及位置类似含四丙基胺模板剂的 ZSM-5 分子筛（TPA-ZSM-5）的 \ce{^{29}Si} 谱。

\textbf{（2）CP 强度的线性区。} 积炭量由 0.85\% 变化到 9.6\% 时，\ce{^{29}Si} CP/MAS NMR 信号强度与积炭量呈线性关系；而积炭量为 23.6\% 时，\ce{^{29}Si} CP 信号强度只是积炭 9.6\% 时的一半，原因是高积炭分子筛表面上生成大量石墨，\ce{^{1}H} 数目减少，降低了 \ce{^{1}H}$\rightarrow$\ce{^{29}Si} 交叉极化效率。

\textbf{（3）H/C 值随积炭量下降。} 在 HZSM-5 上测得：积炭量 0.8\% 时 H/C 为 0.9；积炭量 9.6\% 时 H/C 为 0.5；积炭量 23.6\% 时 H/C 为 0.1。氢碳比从接近 1 降到 0.1，与 \ce{^{13}C} 谱上芳构化的图像一致，说明积炭在高积炭量时已高度脱氢。
\end{dataflow}

\paragraph{（2）\ce{^{27}Al} MAS NMR：脱铝、对称性与「NMR 不可见」的铝}

\begin{principle}
\ce{^{27}Al} 自旋 $I=5/2$，是四极核，谱线主要由四极相互作用控制：一阶四极作用由 MAS 消除，中心跃迁剩下的二阶增宽不能由 MAS 消除，其线宽近似随 $\Delta\nu_{1/2}\propto C_Q^{2}/\nu_0$ 变化，其中 $C_Q=e^{2}qQ/h$ 为四极偶合常数，$\nu_0$ 为共振频率。谱线还受自旋--自旋弛豫限制，$\Delta\nu_{1/2}\approx 1/(\pi T_2)$。积炭后谱的变化是这两种增宽机制共同作用的结果：积炭分子取代了分子筛中的水使水含量减少，而且使 Brönsted 酸位上的 H$^{+}$ 在铝核周围停留更长的时间，从而改变了铝核周围的电场梯度，引起四极相互作用的加强和位移；随着 \ce{^{27}Al} 峰宽的增加，骨架 \ce{^{27}Al} 化学位移向高场移动。积炭后 \ce{^{27}Al} 谱的变化与吸附烷烃和催化剂脱水效应很类似。
\tocheck{未对照原页面图像：二阶四极增宽与线宽的表达式按标准形式补入（原书以文字描述），是否含 $(1+\eta^{2}/3)$ 因子及具体系数需核对原页}

\textbf{从谱变化到失活类型的推断链。} 积炭量增大时，积炭分子覆盖了分子筛的 Brönsted 酸位，改变了骨架铝的对称性且增大了 \ce{^{27}Al} 的四极偶合常数；四极偶合常数越大，\ce{^{27}Al} 信号越弱，NMR 可观察到的铝核越少。这种情况下，催化剂失活主要是由 Brönsted 酸位中毒引起的。
\end{principle}

\begin{measure}
\ce{^{29}Si}：循环延迟要按测得的 $T_1$ 设定（本书实例用 \qty{10}{s}），比较积炭前后强度前先确认样品 $T_1$ 是否改变；CP/MAS 需注明接触时间，并用直接极化谱作参照。

\ce{^{27}Al}：四极核建议用高场、小扳转角（短而强的脉冲）与尽可能高的转速，以减小二阶四极增宽与边带；必须报转速、脉冲条件与参考物（骨架四配位铝约 60 ppm、非骨架铝约 0 ppm、畸变晶格铝约 30 ppm 的归属都建立在同一参考标度上）。不同磁场、不同转速下测得的峰强与峰宽\textbf{不可}直接比较。
\end{measure}

\begin{dataflow}
\textbf{（1）三条曲线的走向。} 积炭量对骨架 \ce{^{27}Al} 信号强度（a）、峰宽（b）和化学位移（c）的影响（图 8-18）规律一致：随积炭量增加，骨架 \ce{^{27}Al} 峰强度降低、峰宽增加、化学位移向高场移动。

\textbf{（2）非骨架铝参与积炭。} 图 8-19 是脱铝 HZSM-5 分子筛上丙酮转化期间的 \ce{^{27}Al} MAS NMR 谱。初始样品出现两个共振信号：60 ppm 的 A 信号为四配位骨架 Al，0 ppm 的 R 信号为非骨架 Al。积炭量为 10.7\% 时，A 峰变宽，非骨架铝峰变成一个肩峰；把催化剂部分氧化使积炭量降为 5.5\% 时，A 峰重新窄化，非骨架铝峰重新出现，只是强度比初始状态低。这种现象说明在 HZSM-5 分子筛上，非骨架铝物种参与了积炭的生成，而且这一过程在部分氧化后可以部分逆转。
\tocheck{未对照原页面图像：图 8-19 的峰标签在 OCR 中前后不一致（先记为 A、R，随后又出现 B），此处按正文归属写作 A（骨架铝 60 ppm）与非骨架铝峰（0 ppm）}

\textbf{（3）骨架铝也会被卷进去。} 在脱铝的 HY 分子筛上，正己烷裂解积炭样品的 \ce{^{27}Al} MAS NMR 谱中，除了骨架铝和非骨架铝信号外，还出现了一个 30 ppm 共振峰，归属为易变形和易碎的晶格铝；随积炭量增加，该峰强度增加，说明积炭后引起某些骨架铝破碎，形成了 30 ppm 铝物种。因此在相当高积炭量的分子筛中，一些骨架铝也卷入了积炭的生成。

\textbf{（4）两种核合起来能给出什么。} 综合分析 \ce{^{29}Si} 和 \ce{^{27}Al} MAS NMR 实验结果，可以研究积炭对分子筛的影响、积炭后分子筛的脱水（及相应的骨架环境变化）和积炭占有的体积；结合催化活性的测量，进而得到酸位中毒或孔道堵塞的催化剂失活信息。
\tocheck{未对照原页面图像：原文「积炭后分子筛脱氧和脱水量」中的「脱氧」含义存疑（疑为 OCR 或原书表述），此处按脱水与体积占据理解}
\end{dataflow}

\begin{pitfall}
\textbf{（1）\ce{^{27}Al} 无信号，不等于没有铝物种。} 四极增宽使部分铝核「NMR 不可见」：积炭覆盖 B 酸位后四极偶合常数增大，\ce{^{27}Al} 信号变弱，可观察到的铝核变少。把「谱上没有这个峰」直接读成「样品里没有这种铝」，是本节最容易犯的错。提高磁场有助于观察，但不同磁场下的强度不可相互换算。

\textbf{（2）\ce{^{27}Al} 强度下降不等于脱铝量。} 强度同时受 $C_Q$ 与 $T_2$ 影响，要判断是否真的脱铝，需与 \ce{^{29}Si} 的 Q$^n$ 分布或化学分析（如 ICP）对照。

\textbf{（3）30 ppm 峰不能直接当作「新的骨架铝物种」的浓度。} 它的归属是易变形、易碎的晶格铝，其强度依赖线型分解的可靠性；在只有一条宽包络时，把它分出来本身就带主观性。

\textbf{（4）\ce{^{29}Si} 信号下降不等于骨架硅减少。} 积炭使 Si 的 $T_1$ 变长，循环时间不足就会造成「假损失」；比较积炭前后必须先测 $T_1$ 或统一使用足够长的循环延迟。同理，\ce{^{29}Si} CP 谱的强度受 \ce{^{1}H} 数目控制，在高积炭量时失效（23.6\% 时只有 9.6\% 时的一半）。
\end{pitfall}

\begin{crosscheck}
\ce{^{29}Si} 与 \ce{^{27}Al} 是一对骨架报告者，但两者都可能因非化学原因变弱（前者因 $T_1$ 变长，后者因 $C_Q$ 增大）。因此「骨架损伤」的结论至少要与 \ce{^{1}H} 的 B 酸位计数、催化活性测量三方对照：单靠 \ce{^{27}Al} 强度下降，既可能是脱铝，也可能是四极增宽把铝藏了起来；而孔道的几何后果（窗口是否堵死、笼是否缩小）由 \ce{^{129}Xe} 回答（§8.6.4）。
\end{crosscheck}

\subsubsection{\ce{^{1}H} MAS NMR 研究分子筛积炭}
\src{S1 §8.6.3，书内 p.559--560 / PDF p.51--52}

\ce{^{1}H} MAS NMR 也是研究分子筛积炭的有效方法：它可以提供积炭后分子筛中仍有活性的 Brönsted 酸位信息。通过脉冲场梯度的自扩散测量可以确定积炭的位置和孔内与外表面扩散的变化；当已知积炭量时，可通过 \ce{^{1}H} MAS NMR 谱确定 H/C 值，进而确定分子筛晶体孔道内和外表面上积炭化合物的性质；对分子筛中吸附分子进行 \ce{^{1}H} 自旋--晶格和自旋--自旋弛豫时间 $T_1$、$T_2$ 的测量，可以得到积炭分子在分子筛孔道中的分布信息。

\begin{measure}
\textbf{转速是前提}：\ce{^{1}H} 同核偶极增宽很大，转速不足时 2.0 ppm 的 Si--OH 峰与 4.2 ppm 的 B 酸峰无法分辨。必须报转速。

\textbf{定量}：单脉冲谱，循环延迟取 $5T_1$ 以上（或先用反转恢复／饱和恢复测 $T_1$）；\ce{^{1}H} 的背景（探头、空转子、吸附水）要单独采集并扣除。峰位归属：2.0 ppm 为 Si--OH，4.2 ppm 为 B 酸中心。

\textbf{弛豫与自扩散}：$T_1$ 用反转恢复或饱和恢复；$T_2$ 用回波序列测量（固体中直接读 FID 衰减会低估 $T_2$）；脉冲场梯度（PFG）\ce{^{1}H} NMR 给出孔内与晶间的自扩散系数，用于判断积炭位置与扩散通道的变化。
\end{measure}

\begin{dataflow}
\textbf{（1）两种探针分子给出两种位置信息。} Lechert 等研究了不同脱铝程度的 HY 分子筛在 530 K 丁烯转化过程中的积炭情况，比较分子筛积炭前后吸附正丁烷和苯的能力及 \ce{^{1}H} 的 $T_1$、$T_2$ 弛豫时间。与积炭前样品相比，积炭后分子筛吸附正丁烷的能力降低不大，但 \ce{^{1}H} 的 $T_1$ 和 $T_2$ 值减小，表明吸附的正丁烷分子的重新取向运动受到很大阻碍，积炭附近分子的堆积发生了变化。吸附苯的情况恰恰相反：\ce{^{1}H} 的 $T_1$ 和 $T_2$ 值变化不大，表明积炭对苯分子的运动影响很小，但吸附能力却大大降低。

\textbf{（2）观察与推断要分开写。} 观察：吸苯量大幅下降，而苯的 $T_1$、$T_2$ 几乎不变。推断：苯分子主要位于分子筛孔道之间的窗口，分子绕着它们的二重轴运动，同时伴随着相邻位之间的跳跃。这一推断解释了「进不去但动得动」的组合，也说明吸附量的变化与弛豫的变化必须分别解读：吸正丁烷是「动不了但进得去」，吸苯是「进不去但动得了」。

\textbf{（3）B 酸位的存活数量。} Ernst 等研究了 HZSM-5 分子筛积炭前后的 \ce{^{1}H} MAS NMR 谱：2.0 ppm 是 Si--OH 峰，4.2 ppm 是 B 酸中心；积炭前 B 酸相对含量很大，积炭后 B 酸的峰强降低，而 Si--OH 含量大大增加，并且出现一个很大的宽峰，这是由积炭分子、或由酸羟基与积炭分子强偶极相互作用所致。根据 B 酸峰强度的降低可以计算出积炭后仍存在的活性 B 酸位的数量；参考 \ce{^{27}Al} MAS NMR 谱中骨架铝积炭后的变化，可分析出积炭使 B 酸位发生了中毒。

\textbf{（4）这一节的出口。} \ce{^{1}H} MAS NMR 谱主要用于估算 B 酸位的数量，结合 \ce{^{27}Al} MAS NMR 谱等确定积炭引起分子筛失活的类型；测量吸附探针分子的弛豫时间和脉冲梯度 \ce{^{1}H} NMR 谱可以提供积炭在分子筛晶体孔道中的分布信息；当已知积炭量时，\ce{^{1}H} MAS NMR 谱可以测量积炭前后的 H/C 值，但要注意：\ce{^{1}H} MAS NMR 谱测出的 H/C 值是含孔道内和外表面全部积炭的数值。
\end{dataflow}

\begin{pitfall}
\textbf{（1）两个 H/C 的口径不同。} 由 \ce{^{29}Si} CP 强度所得的 H/C 只在 CP 强度与积炭量成线性的区间可靠（本例为 0.85\%--9.6\%），高积炭时表面石墨使 \ce{^{1}H} 减少而系统性偏低；由 \ce{^{1}H} MAS NMR 谱所得的 H/C 包含孔道内与外表面的全部积炭。报 H/C 时必须同时说明口径与积炭量区间，两个数值不一致并不矛盾。

\textbf{（2）宽峰不是新酸位。} 积炭后出现的很大宽峰可由积炭分子或酸羟基与积炭分子的强偶极相互作用造成，属于相互作用导致的增宽；把它当成「新增的酸性物种」会同时高估酸位数量。

\textbf{（3）分峰不当会双向失真。} 该宽峰与 4.2 ppm 的 B 酸峰叠加，分峰不当会同时抬高「剩余 B 酸」与「新增物种」；报剩余 B 酸比例时应给出分峰模型与误差。

\textbf{（4）背景水与 Si--OH 重叠。} 2.0 ppm 附近的 Si--OH 与探头／吸附水的 \ce{^{1}H} 背景重叠，判断「Si--OH 大大增加」之前必须扣背景，否则会把水的变化读成羟基的变化。
\end{pitfall}

\begin{crosscheck}
\ce{^{1}H} 给出的是「还剩多少 B 酸位」，\ce{^{27}Al} 给出的是「骨架铝变成了什么」。只有两者对照，才能把 B 酸峰的降低归因于积炭覆盖／中毒，而不是骨架脱铝本身；再叠加活性数据，才能把「酸位中毒」与「孔道堵塞」两种失活类型分开。
\end{crosscheck}

\subsubsection{吸附氙的 \ce{^{129}Xe} NMR 研究分子筛积炭}
\src{S1 §8.6.4，书内 p.560--562 / PDF p.52--54}

测定不同分子筛在不同反应条件下积炭后的吸附氙 $\delta=f(n)$ 曲线，可获得分子筛结构和反应物性质对积炭影响的信息，确定积炭性质和积炭后分子筛结构的变化。Xe 是惰性原子，不参与反应、不改变酸性，它的化学位移由所处空间的尺寸与碰撞环境决定，因此对孔道尺寸与堵塞极其敏感。在积炭研究中，它是少数能直接反映\textbf{孔道空间占据}几何的探针：\ce{^{13}C} 告诉你积炭是什么，\ce{^{129}Xe} 告诉你孔道被占成了什么形状。

\begin{principle}
在快速交换极限下，测得的 \ce{^{129}Xe} 化学位移与 Xe 吸附量 $n$ 的关系可写成线性形式
\[
\delta(n)=\delta_{n\to 0}+k\,n
\]
其中 $\delta_{n\to 0}$ 是外推到零氙吸附量时的位移，反映 Xe 与孔壁及强吸附位（例如非骨架铝）的相互作用，以及在孔口受阻时扩散受限的程度；斜率 $k$ 反映 Xe--Xe 碰撞的频率，与 Xe 可自由活动的体积成反比。由此得到三条判读规则：

（1）两曲线平行（$k$ 相同）而 $\delta_{n\to 0}$ 增大：自由体积基本不变，但 Xe 在笼间的扩散受阻，积炭主要落在窗口（笼与笼之间的通道）上。

（2）$k$ 与 $\delta_{n\to 0}$ 同时增大：自由体积本身下降，且有部分积炭位于晶体外表面，阻碍 Xe 进入。

（3）由积炭与未积炭样品的斜率之比，可以估算可吸附 Xe 的体积下降幅度：
\[
\frac{V_{\text{积炭样品}}}{V_{\text{未积炭样品}}}\approx\frac{k_{\text{未积炭样品}}}{k_{\text{积炭样品}}}
\]
本节实例给出：丙烯裂解生成 15\% 与 33\% 积炭时，可吸附 Xe 的体积分别降低 52\% 与 90\%。

要记住两条轴的物理含义：\textbf{斜率的变化说明自由体积变了，斜率不变而截距变化说明扩散或强吸附位变了}。例如 15\% 积炭样品的斜率与 10.5\% 积炭样品相同，说明此时积炭主要影响晶孔的外表面。
\tocheck{未对照原页面图像：$\delta=f(n)$ 的线性表达、斜率比与可吸附体积下降的换算式（是否取倒数、是否需要独立的 Xe--Xe 参考位移标定）原书只给结果未给公式，此处按常用形式重建}
\tocheck{未对照原页面图像：\ce{^{129}Xe} 天然丰度与「Xe 原子的动态直径约 0.44 nm」为按标准值补入，原书 §8.6.4 只写「减少到约为 Xe 原子的动态直径」而未给数值}
\end{principle}

\begin{measure}
\textbf{样品}：密封管／转子中定量引入 Xe（由压力与体积标定吸附量 $n$），并等待吸附平衡；同一批样品必须做未积炭参照曲线。

\textbf{变温是必需项}：本节的判读大量依赖「斜率与 $\delta_{n\to 0}$ 是否随温度变化」，因此必须做温度序列；只测一条等温曲线，无法区分「自由体积下降」与「扩散受阻」。

\textbf{参考与延迟}：以氙气共振峰为 0 ppm（本节实例中 0 ppm 峰即氙气峰）；\ce{^{129}Xe} 的 $T_1$ 较长，循环延迟要足够。Xe 化学位移跨度大，对磁场要求不高，但对 Xe 量的标定要求高；比较时用斜率与截距，而不是某个单压力点的位移值。
\end{measure}

\begin{dataflow}
\textbf{（1）窗口堵塞的判据。} 1988 年 Ito 等首先将 \ce{^{129}Xe} NMR 用于正己烷及丙烯裂化的 HY 分子筛积炭研究。正己烷为反应物时，未积炭样品（HY）与积炭量 5.5\%（He-5）的 $\delta=f(n)$ 曲线互相平行，说明二者斜率相同，但 $\delta_{n\to 0}$ 值不同；He-5 的 $\delta_{n\to 0}$ 增大，表明氙在超笼间的扩散减少，积炭位于超笼之间的窗口上。

\textbf{（2）自由体积下降与外表面堵塞。} 丙烯为反应物、积炭分别为 15\%（Pr-15）和 33\%（Pr-33）时，与未积炭样品相比，$\delta=f(n)$ 直线的斜率及 $\delta_{n\to 0}$ 均增加很大；根据积炭和未积炭样品斜率之比，可计算出积炭 15\% 和 33\% 的分子筛可吸附 Xe 的体积分别降低 52\% 和 90\%。这表明积炭后不仅分子筛孔体积大大降低，而且部分积炭位于晶体的外表面，阻碍 Xe 进入超笼。33\% 积炭样品的 \ce{^{129}Xe} 谱出现两个峰：一个是很宽且低场位移很大的谱峰（157 ppm），为吸附在与积炭相连的超笼中的 \ce{^{129}Xe}；另一峰出现在 10 ppm，与 Xe 压力无关，是吸附在分子筛孔道内积炭形成的中孔或微笼中的 \ce{^{129}Xe}。这些结果直接显示了积炭的位置。

\textbf{（3）初始积炭落在非骨架铝附近。} Barrage 等系统研究了庚烷裂解时失活的 HY 分子筛中积炭的分布。未积炭样品的 $\delta=f(n)$ 曲线在低 Xe 浓度时出现一个最低点，并且随温度升高 $\delta_{n\to 0}$ 增大，这是 Xe 与非骨架铝的强吸附位相互作用的特征，\ce{^{27}Al} NMR 谱也证明了这一点。当积炭量为 3\%、10.5\% 和 15\% 时，$\delta=f(n)$ 曲线全部变成直线，这表明最初始的积炭位于或接近于非骨架 Al 物种，积炭阻碍了非骨架 Al 物种与 Xe 原子的相互作用，消除了曲线中的最低点。而且，当积炭量达到 3\% 时，庚烷裂解的活性丧失一半多，这说明非骨架 Al 物种在反应初始对 HY 分子筛裂解活性起主导作用。

\textbf{（4）积炭由非均匀走向均匀的三段图像。}
低积炭量（3\%）：斜率变化较小、与初始样品相近，自由体积几乎不受影响，但 $\delta_{n\to 0}$ 增加了，且两者随温度升高而发生变化，说明 Xe 的扩散受阻，积炭初始是非均匀的并且主要沉积在超笼的窗口上。
中等积炭量（10\%）：$\delta=f(n)$ 曲线的斜率和 $\delta_{n\to 0}$ 都增加，但与吸附温度无关，此时积炭的生成更均匀，且超笼内也生成积炭，将其直径减少到约为 Xe 原子的动态直径；温度升高时 Xe 原子扩散未受到阻碍，所以斜率增加，而 $\delta=f(n)$ 曲线与温度无关。
高积炭量（15\%）：\ce{^{129}Xe} 谱出现三条谱线，低场的一条宽峰，其化学位移值随 Xe 浓度呈线性变化；第二条谱线出现在 44 ppm，与 Xe 压力无关；0 ppm 峰是氙气的共振峰。三条谱线说明 Xe 吸附在不同的位置上，相互之间在 NMR 时域内没有交换。其 $\delta=f(n)$ 曲线斜率与 10.5\% 积炭量的样品相同，但 $\delta$ 值增加很大，表明 Xe 的平均自由通道很小、扩散严重受阻、许多笼被积炭堵塞：积炭不仅大大减小了内部自由体积，而且在分子筛晶孔之间形成了微孔或中孔；斜率相同表明积炭主要影响晶孔的外表面。

\textbf{（5）可推广性。} 以同样的方法，人们对 HZSM-5 分子筛的积炭也进行了大量研究，确定了不同积炭量时积炭的分布。总之，分子筛吸附 Xe 后测定 \ce{^{129}Xe} NMR 的 $\delta=f(n)$ 曲线，可以直接得到内部积炭位置及积炭量信息，能够显示分子筛结构和反应物性质对积炭和再生的影响。
\end{dataflow}

\begin{pitfall}
\textbf{（1）多条谱线不等于多种积炭物种。} 33\% 样品的两个峰（157 ppm 的宽峰、10 ppm 且与压力无关的峰）与 15\% 样品的三条谱线（低场宽峰、44 ppm、0 ppm），说明的是 Xe 处于不同位置且在 NMR 时域内互不交换。其中 0 ppm 峰是氙气本身，不能计入孔内物种；「与压力无关」是判断该峰来自封闭小空间（中孔、微笼）的关键证据。

\textbf{（2）位移变化的来源不止「孔变小」。} 非骨架铝的强吸附位与温度都会改变 $\delta$ 与曲线形状。未积炭样品的曲线最低点消失，是积炭覆盖了强吸附位，而不是孔容恢复；不做未积炭对照、不做变温，很容易把这一现象误读为「积炭改善了孔道」。

\textbf{（3）温度依赖性是判据而不是噪声。} 斜率是否随温度变化、$\delta_{n\to 0}$ 是否随温度变化，分别对应自由体积与扩散受阻；两者必须分别报告，否则 3\% 与 10\% 这两段图像无法区分。

\textbf{（4）Xe 进不去的孔，测不到。} 探针法的固有边界：被完全堵死的空间不产生 \ce{^{129}Xe} 信号，所以 \ce{^{129}Xe} 不会「报告」那部分孔容损失；把斜率比换算成体积下降，也依赖一个独立的体积标定。
\end{pitfall}

\begin{crosscheck}
\textbf{（1）与氮吸附孔容下降的互补关系。} 77 K 氮吸附给出的是孔容／比表面积的总量下降，是一个标量：它能告诉你「少了多少」，但不告诉你积炭落在窗口、笼内还是外表面。\ce{^{129}Xe} 的 $\delta=f(n)$ 恰好补上「位置」这一维：曲线平行而截距升高对应窗口堵塞；斜率与截距同时升高对应笼内自由体积下降并伴以外表面堵塞；出现与压力无关的新峰对应积炭内部新生成了微孔或中孔；$\delta_{n\to 0}$ 的温度依赖对应扩散是否受阻。反过来，\ce{^{129}Xe} 只能报告 Xe 能到达的空间，完全堵死的孔不产生信号，因此「孔容一共降了多少」仍需氮吸附这类总量方法提供基准，而把斜率之比换算成体积下降也依赖独立标定。两者的关系可以一句话概括：\textbf{氮吸附给总量（多少），\ce{^{129}Xe} 给分布（在哪、以什么方式堵）}，二者不可互相替代。（原书 §8.6.4 未与氮吸附数据对照，这一互补关系是按方法学写出的推论。）

\textbf{（2）与其余三种核的合围。} \ce{^{13}C} 给积炭的化学身份，\ce{^{1}H} 给剩余 B 酸位与 H/C，\ce{^{29}Si}／\ce{^{27}Al} 给骨架损伤，\ce{^{129}Xe} 给孔道几何。本节中一个可用的联动判据是：积炭只有 3\% 时裂解活性已丧失一半多，而自由体积几乎不受影响（斜率与初始样品相近），这是典型的酸位中毒而不是孔道堵塞，与 \ce{^{27}Al} 谱中非骨架铝强吸附位被积炭覆盖的观察一致。
\end{crosscheck}

\begin{methodsum}
积炭的四问与四种核的对应关系可以压缩成一句话：\textbf{\ce{^{13}C} 回答「是什么」，\ce{^{29}Si}／\ce{^{27}Al} 回答「骨架被弄坏到什么程度」，\ce{^{1}H} 回答「还剩多少酸位、氢在谁身上」，\ce{^{129}Xe} 回答「堵在哪里、堵成什么样」。}

组合使用有三条纪律：（1）凡涉及「量」的数字，必须说明是直接极化还是 CP、有无独立标定；（2）凡涉及「消失」的结论，必须先排除弛豫变慢、四极增宽、旋转边带与探头调谐等非化学原因；（3）凡涉及「位置」的结论，至少要有两个独立方法支持（例如窗口堵塞同时由 \ce{^{129}Xe} 的截距升高与 \ce{^{1}H} 上 B 酸位存活数佐证）。
\end{methodsum}

\begin{table}[tbp]
\centering
\caption{四种核在分子筛积炭研究中的分工对照。读者应把同一行的四项一起读：「测什么」决定它能回答哪个问题，
「定量前提」决定数字能不能当量用，「主要干扰」决定结论的边界；只取信号强度、不看后两列，是本领域最常见的误用。}
\label{tab:nuclei-coke}
{\srcsize
\begin{tabularx}{\twocolwidth}{@{}p{13mm}p{27mm}p{45mm}p{38mm}X@{}}
\toprule
核种 & 测什么 & 能给出的积炭信息 & 定量前提 & 主要干扰 \\
\midrule
\ce{^{13}C} & 碳的化学环境（0--300 ppm） &
脂肪碳 10--40 ppm、烯碳约 150 ppm、芳香碳 125--145 ppm、聚芳烃宽芳香区 130--170 ppm；随分子筛、反应物与温度改变 &
CP 谱不能定量；若报碳量需直接极化谱 + 足够长循环延迟 + 已知标样（或热重／元素分析） &
各类型碳 CP 效率不同；旋转边带伪装成物种；高积炭时导电石墨与顺磁中心使信号丢失（23.6\% 积炭仅 29\% 可检出） \\
\ce{^{29}Si} & 骨架硅的 Q$^n$ 环境与 $T_1$ &
积炭对骨架的扰动：峰移向低场、半峰宽 \qty{143}{Hz} 增至 \qty{218}{Hz}、信号下降；CP 强度给出 H/C（0.8\%、9.6\%、23.6\% 积炭时为 0.9、0.5、0.1） &
CP 强度按每晶胞含氢数标定；H/C 需同时知道积炭量；只在 CP 强度与积炭量成线性的区间（0.85\%--9.6\%）可靠 &
积炭取代孔道中顺磁性氧使 Si 的 $T_1$ 变长，循环时间不足即出现「假损失」；表面石墨减少 \ce{^{1}H} 使 CP 效率下降 \\
\ce{^{27}Al} & 骨架四配位 Al（约 60 ppm）、非骨架 Al（约 0 ppm）、畸变晶格 Al（约 30 ppm） &
脱铝与骨架损伤；积炭覆盖 B 酸位使骨架铝对称性下降、四极偶合常数增大；非骨架铝与部分骨架铝参与积炭生成 &
需同磁场、同转速、同脉冲条件下比较；峰强同时受 $C_Q$ 与 $T_2$ 影响，不能直接换算成铝量或脱铝量 &
四极增宽使部分铝「NMR 不可见」，无信号不等于无物种；强度下降未必是脱铝；30 ppm 峰依赖线型分解 \\
\ce{^{1}H} & Si--OH（约 2.0 ppm）、B 酸中心（约 4.2 ppm）；探针分子的 $T_1$、$T_2$ 与自扩散 &
积炭后仍存活的 B 酸位数量（据此判断酸位中毒）；H/C 值；积炭在孔道内与外表面的分布（弛豫 + 脉冲场梯度自扩散） &
需已知积炭量；B 酸定量需与 \ce{^{27}Al} 对照，并用 $5T_1$ 以上循环延迟、扣除探头与水背景 &
探头／水背景与 Si--OH 峰重叠；积炭宽峰与 B 酸峰的强偶极重叠使分峰失真；所得 H/C 含内、外表面全部积炭 \\
\ce{^{129}Xe} & 化学位移与 Xe 吸附量的关系 $\delta=f(n)$（斜率与 $\delta_{n\to 0}$） &
积炭位置（超笼窗口／笼内／外表面）；可吸附 Xe 体积下降（15\%、33\% 积炭分别降 52\%、90\%）；谱峰数目给出互不交换的 Xe 环境（如 157 ppm、10 ppm、44 ppm） &
需与未积炭样品比较斜率与截距；变温测量是把「自由体积下降」与「扩散受阻」分开的前提；Xe 量须定量标定 &
温度与强吸附位（非骨架 Al）同样改变 $\delta$；多峰不等于多物种（0 ppm 为氙气）；Xe 进不去的孔完全测不到 \\
\bottomrule
\end{tabularx}}
\end{table}

\srcnote{本节覆盖 S1 §8.6（书内 p.555--562 / PDF p.47--54）：积炭失活的基本图像与四种核的分工；其中 \ce{^{13}C} 部分见 §8.6.1，\ce{^{29}Si}／\ce{^{27}Al} 部分见 §8.6.2，\ce{^{1}H} 部分见 §8.6.3，\ce{^{129}Xe} 部分见 §8.6.4。}
